\chapter{Risultati Sperimentali}

\section{Introduzione}
Per valutare l'efficacia del sistema di analisi proposto, è stata
progettata una suite di benchmark che misura la capacità
dell'analizzatore di estrarre correttamente le entità architetturali
(nodi) e le relazioni di dipendenza (archi) da progetti
multi-linguaggio. Questo capitolo presenta la metodologia di
validazione, i risultati aggregati per linguaggio e un'analisi
qualitativa delle categorie di errore riscontrate.

\section{Metodologia di Validazione}
La validazione si basa sul confronto tra il grafo delle dipendenze
estratto automaticamente dal sistema e un \textit{ground truth}
definito manualmente per ciascun benchmark. Per ogni linguaggio
supportato --- Rust, Java, Python, C e C++ --- è stato costruito un
progetto di test che esercita i principali costrutti architetturali:
definizioni di moduli, tipi strutturati (classi, struct, interfacce,
trait, enum), funzioni e metodi, relazioni di ereditarietà, chiamate,
istanziazioni e accessi a campo.

Il ground truth è codificato in un file di manifesto YAML
(\texttt{test.yml}) che specifica:
\begin{itemize}
  \item \textbf{Nodi attesi}: le entità architetturali che
    l'analizzatore deve riconoscere, con il rispettivo tipo
    (\texttt{module}, \texttt{class}, \texttt{struct}, \texttt{function},
    \texttt{field}, ecc.).
  \item \textbf{Archi attesi}: le dipendenze tra entità, con sorgente,
    destinazione e tipo di relazione (\texttt{uses\_type},
      \texttt{calls}, \texttt{inherits}, \texttt{accesses\_field},
    \texttt{imports}, ecc.).
\end{itemize}

Un eseguibile dedicato (\texttt{benchmark\_runner}) esegue l'analisi
sul progetto di test e confronta il risultato con il manifesto,
generando un report dettagliato sia in formato Markdown (per la
consultazione umana) sia in formato JSON (per l'elaborazione
automatizzata).

\subsection{Metriche}
Per ciascun benchmark, vengono calcolate due metriche fondamentali:
\begin{itemize}
  \item \textbf{Accuratezza Nodi} (\textit{Node Accuracy}): la
    frazione di nodi attesi correttamente individuati dall'analizzatore,
    sia in termini di esistenza che di classificazione del tipo.
    $$\text{Acc}_{nodi} = \frac{|\{n \in N_{attesi} \mid n \in
    N_{estratti} \land \text{kind}(n) \text{ corretto}\}|}{|N_{attesi}|}$$
  \item \textbf{Accuratezza Archi} (\textit{Edge Accuracy}): la
    frazione di archi attesi correttamente estratti, dove un arco è
    considerato corretto se sorgente, destinazione e tipo di relazione
    corrispondono al ground truth.
    $$\text{Acc}_{archi} = \frac{|\{e \in E_{attesi} \mid e \in
    E_{estratti} \land \text{kind}(e) \text{ corretto}\}|}{|E_{attesi}|}$$
\end{itemize}

\section{Risultati Aggregati}
La Tabella~\ref{tab:risultati-aggregati} riporta i risultati
complessivi per ciascun linguaggio.

\begin{table}[ht]
  \centering
  \caption{Risultati dei benchmark per linguaggio.}
  \label{tab:risultati-aggregati}
  \begin{tabular}{l r r r r r r}
    \hline
    \textbf{Linguaggio} & \textbf{Nodi Attesi} & \textbf{Trovati} &
    \textbf{Acc. Nodi} & \textbf{Archi Attesi} & \textbf{Trovati} &
    \textbf{Acc. Archi} \\
    \hline
    Rust   & 39 & 38 & 97{,}4\% & 67 & 46 & 68{,}7\% \\
    Java   & --\textsuperscript{*} & -- & -- & 34 & 28 & 82{,}4\% \\
    Python & 16 & 16 & 100\%    & 19 & 12 & 63{,}2\% \\
    C      &  7 &  7 & 100\%    & 10 &  4 & 40{,}0\% \\
    C++    & 12 & 12 & 100\%    & 11 &  2 & 18{,}2\% \\
    \hline
  \end{tabular}
  \begin{flushleft}
    \textsuperscript{*}\footnotesize{Il benchmark Java è basato su un
      progetto architetturale reale (HUSACCT) e testa esclusivamente la
    correttezza degli archi di dipendenza cross-package.}
  \end{flushleft}
\end{table}

\subsection{Osservazioni Generali}
I risultati evidenziano una chiara distinzione tra la capacità del
sistema di riconoscere le entità architetturali e quella di ricostruire
le relazioni di dipendenza tra di esse.

L'\textbf{accuratezza sui nodi} è molto alta in tutti i linguaggi
testati: 100\% per Python, C e C++, e 97{,}4\% per Rust. L'unico
nodo mancante in Rust è una variabile statica a livello di modulo
(\texttt{functions.STATIC\_SA}), il cui mancato riconoscimento è
legato a una limitazione nell'euristica di estrazione delle variabili
dichiarate all'interno di blocchi macro. Questo risultato conferma che
le euristiche di classificazione dei nodi CST nell'IR ---
indipendentemente dal linguaggio sorgente --- sono efficaci e robuste.

L'\textbf{accuratezza sugli archi} presenta invece una variabilità
significativa tra linguaggi. Java raggiunge l'82{,}4\%, Rust il
68{,}7\%, Python il 63{,}2\%, mentre C e C++ ottengono
rispettivamente il 40{,}0\% e il 18{,}2\%. Questa variabilità non è
tuttavia distribuita uniformemente tra le categorie di dipendenza: come
discusso nella Sezione~\ref{sec:analisi-errori}, gli errori sono
concentrati in un numero ristretto di scenari ben identificati.

\section{Analisi per Categoria di Dipendenza}
\label{sec:analisi-errori}
Per comprendere la distribuzione degli errori, i risultati vengono
disaggregati per categoria di arco. La
Tabella~\ref{tab:risultati-categoria} riporta l'accuratezza per
ciascuna tipologia di dipendenza.

\begin{table}[ht]
  \centering
  \caption{Accuratezza per categoria di dipendenza (archi).}
  \label{tab:risultati-categoria}
  \begin{tabular}{l c c c c c}
    \hline
    \textbf{Categoria} & \textbf{Rust} & \textbf{Java} &
    \textbf{Python} & \textbf{C} & \textbf{C++} \\
    \hline
    Import / Include       & 14/14 & --\textsuperscript{a}    & 3/3
    & --    & --   \\
    Ereditarietà (IsA)     & 1/1   & 3/3  & 1/1   & --    & 1/1  \\
    Uses Type (dichiaraz.) & 11/18 & 15/16 & --   & 1/2   & 1/3  \\
    Calls (invocazioni)    & 8/11  & 8/8  & 3/7   & 3/4   & 0/3  \\
    Accesses Field         & 4/4   & 11/11 & 6/8  & 0/4   & 0/4  \\
    Type Cast              & 0/1   & 0/2  & --    & --    & --   \\
    \hline
  \end{tabular}
  \begin{flushleft}
    \textsuperscript{a}\footnotesize{Il benchmark Java non testa gli
      import come archi espliciti nel ground truth, ma gli import sono
    necessariamente risolti per poter estrarre le altre dipendenze.}
  \end{flushleft}
\end{table}

\subsection{Categorie con Alta Accuratezza}

\paragraph{Import e Risoluzione dei Moduli.}
La risoluzione degli import è pienamente funzionante in tutti i
linguaggi testati: Rust (14/14), Python (3/3). In Java, gli import
non sono testati esplicitamente come archi nel ground truth, ma il
funzionamento corretto della risoluzione degli archi cross-package
(come l'accesso a \texttt{technology.direct.dao.ProfileDAO} dal
package \texttt{domain.direct.violating}) conferma implicitamente il
corretto funzionamento del sistema di import.

\paragraph{Ereditarietà.}
Le relazioni di ereditarietà (\texttt{IsA}, \texttt{isChildOf},
\texttt{isImplementationOf}) sono correttamente estratte in tutti i
linguaggi. In Java, il benchmark verifica con successo sia
l'ereditarietà da classi concrete, sia l'estensione di classi
astratte, sia l'implementazione di interfacce. In C++, l'ereditarietà
da \texttt{Transport::Vehicle} a \texttt{Transport::Car} è
correttamente identificata. In Rust, la relazione tra trait
(\texttt{TraitB : TraitA}) è riconosciuta tramite il meccanismo di
\texttt{super\_types}.

\paragraph{Invocazioni di Metodo (Java).}
Il benchmark Java mostra un'accuratezza del 100\% sulle chiamate a
metodo (8/8), incluse chiamate a metodi di istanza, metodi statici,
metodi su superclassi, metodi su inner class, e persino chiamate a
librerie esterne (\texttt{FoursquareApi.getAuthenticationUrl}). Questo
risultato dimostra la maturità del Query Engine Algebrico nel contesto
di un linguaggio con tipizzazione statica esplicita.

\paragraph{Accesso ai Campi (Java e Python, contesto d'istanza).}
L'accesso ai campi all'interno di metodi d'istanza funziona
correttamente sia in Java (11/11, inclusi accessi a campi di
superclassi e costanti di interfaccia) sia in Python (6/8,
all'interno dei metodi \texttt{\_\_init\_\_} e dei metodi d'istanza).
Questo è reso possibile dall'iniezione anticipata di
\texttt{self}/\texttt{this} nello \texttt{SymbolStack} durante la
Fase 2a.

\subsection{Categorie con Errori Sistematici}

\paragraph{Accesso ai Campi in Funzioni Libere (C e C++).}
L'accuratezza sugli \texttt{accesses\_field} cala drasticamente
quando il contesto è una funzione libera (non un metodo d'istanza): in
C si ottiene 0/4 e in C++ 0/4. La causa è strutturale: nelle funzioni
libere, il tipo della variabile locale su cui avviene l'accesso al
campo (ad esempio, \texttt{rect.width} dove \texttt{rect} è di tipo
\texttt{Rectangle}) deve essere inferito dall'assegnamento o dal
passaggio per parametro. Il sistema attuale risolve il tipo solo
quando è annotato esplicitamente nella dichiarazione della variabile,
ma nei linguaggi C e C++ le variabili locali di tipo struct sono
frequentemente dichiarate senza annotazione esplicita in contesti dove
il tipo è ovvio dal costruttore o dall'inizializzazione.

\paragraph{Campi delle Varianti Enum con Dati Associati (Rust).}
In Rust, le varianti di enum con dati associati (ad esempio
\texttt{EnumB::FIRST(EnumA)}) producono campi nell'IR che non vengono
correttamente collegati al tipo del dato associato. Questo genera
falsi negativi negli archi \texttt{uses\_type} dai campi delle
varianti enum ai tipi delle varianti stesse. Si tratta di una
limitazione nell'euristica di estrazione che non gestisce ancora il
parsing ricorsivo dei campi posizionali (tuple variant) dell'enum di
Rust.

\paragraph{Chiamate via \texttt{super()} in Python.}
Il benchmark Python evidenzia l'incapacità attuale di risolvere le
chiamate a metodi di superclasse tramite \texttt{super()}. Le chiamate
\texttt{super().\_\_init\_\_()} e \texttt{super().get\_info()} non
vengono estratte perché la risoluzione di \texttt{super()} come
espressione che restituisce il tipo della superclasse non è ancora
implementata nel Query Engine.

\paragraph{Type Cast.}
I cast di tipo (\texttt{as} in Rust, \texttt{(Type)} in Java) non
sono attualmente estratti come archi di dipendenza. Il sistema non
riconosce le espressioni di cast come siti di riferimento a tipi, il
che porta a 0/1 in Rust e 0/2 in Java per questa categoria.

\paragraph{Risoluzione Cross-Modulo delle Variabili Locali (Python).}
In Python, due archi \texttt{accesses\_field} nella funzione
\texttt{main.main} (accessi a \texttt{User.username} e
\texttt{User.birth\_year}) non sono risolti. La causa è la mancata
inferenza del tipo di una variabile locale assegnata tramite
costruttore importato: il sistema non collega
\texttt{admin = Admin(...)} al tipo \texttt{models.Admin} perché la
risoluzione delle istanziazioni come tipo non propaga l'informazione
di tipo alle variabili locali nel blocco corrente.

\section{Discussione}
I risultati complessivi mostrano che il sistema raggiunge il suo
obiettivo primario: l'estrazione language-agnostic delle entità
architetturali è efficace e pressoché completa. La pipeline formale
$\Phi = G \circ R \circ E$ produce nodi corretti con un'accuratezza
superiore al 97\% in tutti i linguaggi testati.

Le limitazioni riscontrate negli archi sono concentrate in tre aree
tecniche ben definite:
\begin{enumerate}
  \item \textbf{Inferenza di tipo per variabili locali}: l'assenza di
    un meccanismo di propagazione del tipo dal lato destro
    dell'assegnamento (costruttore, chiamata di funzione) alle
    variabili locali impedisce la risoluzione degli accessi a campo e
    delle chiamate a metodo su variabili il cui tipo non è annotato
    esplicitamente.
  \item \textbf{Costrutti language-specific non ancora mappati}: cast
    di tipo, \texttt{super()}, e campi posizionali delle enum variant
    richiedono estensioni puntuali alle euristiche di estrazione.
  \item \textbf{Campi d'istanza dichiarati in classe}: il supporto
    dei campi come variabili d'istanza (ad esempio
    \texttt{DeclarationVariableInstance} in Java) richiede il
    riconoscimento delle dichiarazioni a livello di corpo di classe
    oltre che a livello di firma.
\end{enumerate}

È importante osservare che nessuna delle limitazioni riscontrate mette
in discussione la validità dell'approccio language-agnostic in sé: le
correzioni necessarie sono localizzate nelle euristiche di estrazione
(Fase 1) e nell'estensione del motore di inferenza (Fase 2b), senza
richiedere modifiche all'architettura complessiva della pipeline o
all'algebra delle query.

\subsection{Confronto Qualitativo con Strumenti Esistenti}
Il benchmark Java utilizzato è basato sul dataset di HUSACCT
(\textit{Hughes Software Architecture Compliance Checking Tool}), uno
strumento di analisi architetturale language-specific per Java. Il
nostro sistema raggiunge un'accuratezza dell'82{,}4\% sugli archi
dello stesso dataset, utilizzando una pipeline completamente
generica. Le 6 dipendenze non rilevate corrispondono a funzionalità
non ancora implementate (cast di tipo, annotazioni come dipendenze, e
campi d'istanza dichiarati nel corpo della classe), non a difetti
nell'architettura del risolutore.

\section{Riepilogo}
La Tabella~\ref{tab:riepilogo} sintetizza lo stato di supporto per
ciascuna categoria di dipendenza nei linguaggi testati.

\begin{table}[ht]
  \centering
  \caption{Stato di supporto per categoria di dipendenza.}
  \label{tab:riepilogo}
  \begin{tabular}{l c c c c c}
    \hline
    \textbf{Funzionalità}        & \textbf{Rust} & \textbf{Java} &
    \textbf{Python} & \textbf{C} & \textbf{C++} \\
    \hline
    Estrazione Entità             & \checkmark & \checkmark &
    \checkmark & \checkmark & \checkmark \\
    Import / Use / Include        & \checkmark & \checkmark &
    \checkmark & --         & --         \\
    Ereditarietà                  & \checkmark & \checkmark &
    \checkmark & --         & \checkmark \\
    Tipo Parametri e Ritorno      & \checkmark & \checkmark &
    Parziale   & \checkmark & Parziale   \\
    Chiamate Metodo (d'istanza)   & \checkmark & \checkmark &
    Parziale   & \checkmark & --         \\
    Accesso Campi (d'istanza)     & \checkmark & \checkmark &
    \checkmark & --         & --         \\
    Accesso Campi (funz. libere)  & --         & --         & --
    & --         & --         \\
    Type Cast                     & --         & --         & --
    & --         & --         \\
    \hline
  \end{tabular}
\end{table}

I risultati ottenuti confermano la validità dell'approccio
language-agnostic e identificano con precisione le aree di
miglioramento, che saranno oggetto di sviluppo nel capitolo dedicato
agli Sviluppi Futuri.
